\chapter{Fourier Characters and Algebraic Normalization}
\label{app:fourier_characters}

\section{Introduction and Scope}

The discrete Fourier transform on a finite lattice $\Lambda = \mathbb{Z}/L\mathbb{Z}$ forms the kinematic bridge between position-space fermions and momentum-space modes. In Chapter~\ref{ch:fourier}, the Fourier transform was formulated over the complex field $\mathbb{C}$ using the canonical primitive root $\zeta = e^{2\pi i / L}$.

However, from an algebraic and foundational perspective, several critical questions arise:
\begin{enumerate}
    \item \emph{What is the minimal algebraic structure required for discrete Fourier duality?} Does character orthogonality hold over general commutative rings, or are integral domains strictly necessary?
    \item \emph{How do we rigorously decouple the algebraic character identities from the metric normalization $1/\sqrt{L}$?} In pure algebra, square roots of integers may not exist, yet exact two-sided Fourier inversion remains valid whenever $L$ is invertible.
    \item \emph{How is representative independence reconciled with the centered momentum band $\Lambda^*$?} The momentum band $\Lambda^* = \{-h+1, \dots, h\}$ is a set of signed integers, not a group under ordinary integer addition.
\end{enumerate}

This appendix provides the complete algebraic answers to these questions. We develop the three-tiered character architecture, prove geometric sum cancellation over arbitrary fields, establish the unscaled operator identities $T S = L I$ and $S T = L I$, and compute the exact lattice dispersion algebraically without relying on transcendental trigonometric functions.

\section{The Minimal Ring Constraint and a Zero-Divisor Counterexample}

It is tempting to develop discrete Fourier analysis over an arbitrary commutative ring $R$ equipped with an element $\zeta \in R$ of multiplicative order $L$. However, this generality fails in the presence of zero divisors.

\begin{example}[Failure of Character Orthogonality over Zero Divisors]
Consider the commutative ring $R = \mathbb{Z}/15\mathbb{Z}$ with lattice size $L = 2$.
The element $\zeta = 4 \in \mathbb{Z}/15\mathbb{Z}$ satisfies:
\begin{equation}
\zeta^1 = 4 \not\equiv 1 \pmod{15}, \qquad \zeta^2 = 16 \equiv 1 \pmod{15}.
\end{equation}
Thus $\zeta = 4$ has exact multiplicative order $L = 2$. Furthermore, the lattice size $L = 2$ is an invertible unit in $R$, since $2 \times 8 = 16 \equiv 1 \pmod{15}$.
However, evaluating the geometric character sum for the non-trivial frequency $k = 1$ yields:
\begin{equation}
\sum_{x=0}^{1} \zeta^{1 \cdot x} = \zeta^0 + \zeta^1 = 1 + 4 = 5 \not\equiv 0 \pmod{15}.
\end{equation}
Even though $L$ is invertible and $\zeta$ has order $L$, the character sum fails to vanish because $(\zeta - 1) = 3$ is a zero divisor in $\mathbb{Z}/15\mathbb{Z}$, and $3 \times 5 = 15 \equiv 0$.
\end{example}

Consequently, character orthogonality cannot be established over general commutative rings. Throughout this appendix, we work over an arbitrary \emph{field} $K$ (or an integral domain) of characteristic not dividing $L$, containing a primitive $L$-th root of unity $\zeta$.

\section{The Three-Tier Character Architecture}

To handle the signed momentum band $\Lambda^* = \{-h+1, \dots, h\}$ without imposing an artificial group law, we define Fourier characters across three distinct algebraic tiers.

\begin{definition}[Tier 1: Integer Character Pairing]
Let $K$ be a field and let $\zeta \in K^\times$ be an element of multiplicative order $L$. The \emph{integer character pairing} is the map $\chi_{\mathbb{Z}}: \mathbb{Z} \times \mathbb{Z} \to K^\times$ defined by:
\begin{equation}
\chi_{\mathbb{Z}}(\zeta, k, x) \coloneqq \zeta^{k x}.
\end{equation}
Because $\zeta \neq 0$, the power $\zeta^{k x}$ is well-defined for all signed integers $k, x \in \mathbb{Z}$ via field inversion $\zeta^{-1}$ for negative exponents.
\end{definition}

\begin{lemma}[Representative Independence]
\label{lem:a01_rep_independence}
Let $\zeta \in K$ be a primitive $L$-th root of unity (so $\zeta^L = 1$). For all integers $k, x, s, t \in \mathbb{Z}$:
\begin{equation}
\chi_{\mathbb{Z}}(\zeta, k + s L, \; x + t L) = \chi_{\mathbb{Z}}(\zeta, k, x).
\end{equation}
\end{lemma}

\begin{proof}
Expanding the exponent of $\zeta$:
\begin{align*}
(k + s L)(x + t L) &= k x + k t L + s x L + s t L^2 \\
&= k x + L (k t + s x + s t L).
\end{align*}
Since $\zeta \neq 0$, exponentiation respects integer addition:
\[
\zeta^{(k + s L)(x + t L)} = \zeta^{k x} \left(\zeta^L\right)^{k t + s x + s t L} = \zeta^{k x} (1)^{k t + s x + s t L} = \zeta^{k x}.
\]
\end{proof}

\begin{definition}[Tier 2: Residue Character]
For residue classes $k, x \in \mathbb{Z}/L\mathbb{Z}$ with canonical unsigned representatives $k_{\mathrm{val}}, x_{\mathrm{val}} \in \{0, 1, \dots, L-1\}$, the \emph{residue character} is:
\begin{equation}
\chi_{\mathrm{res}}(L, \zeta, k, x) \coloneqq \chi_{\mathbb{Z}}(\zeta, k_{\mathrm{val}}, x_{\mathrm{val}}).
\end{equation}
By Lemma~\ref{lem:a01_rep_independence}, $\chi_{\mathrm{res}}(L, \zeta, k, x)$ is independent of the choice of integer representatives for $k$ and $x$.
\end{definition}

\begin{definition}[Tier 3: Band Character]
For a signed momentum mode $k \in \Lambda^*$ and a lattice site $x \in \Lambda = \mathbb{Z}/L\mathbb{Z}$, the \emph{band character} is:
\begin{equation}
\chi_{\mathrm{band}}(L, \zeta, k, x) \coloneqq \chi_{\mathbb{Z}}(\zeta, k, x_{\mathrm{val}}) = \chi_{\mathrm{res}}(L, \zeta, \pi(k), x),
\end{equation}
where $\pi: \Lambda^* \to \mathbb{Z}/L\mathbb{Z}$ is the canonical projection $\pi(k) = k \bmod L$.
\end{definition}

\begin{theorem}[Character Homomorphism Laws]
\label{thm:character_homomorphisms}
Let $\zeta \in K$ be a primitive $L$-th root of unity. The character pairing satisfies:
\begin{enumerate}
    \item \textbf{Spatial Homomorphism:} For all $k \in \mathbb{Z}/L\mathbb{Z}$ and $x, y \in \mathbb{Z}/L\mathbb{Z}$,
    \begin{equation}
    \chi_{\mathrm{res}}(k, x + y) = \chi_{\mathrm{res}}(k, x) \chi_{\mathrm{res}}(k, y).
    \end{equation}
    \item \textbf{Momentum Homomorphism:} For all $k, p \in \mathbb{Z}/L\mathbb{Z}$ and $x \in \mathbb{Z}/L\mathbb{Z}$,
    \begin{equation}
    \chi_{\mathrm{res}}(k + p, x) = \chi_{\mathrm{res}}(k, x) \chi_{\mathrm{res}}(p, x).
    \end{equation}
    \item \textbf{Inversion:} $\chi_{\mathrm{res}}(-k, x) = \chi_{\mathrm{res}}(k, x)^{-1} = \chi_{\mathrm{res}}(k, -x)$.
    \item \textbf{Transported Band Addition:} For $k, p \in \Lambda^*$, let $k \oplus p \in \Lambda^*$ denote the transported band addition defined in Chapter~\ref{ch:lattice_band}. Then:
    \begin{equation}
    \chi_{\mathrm{band}}(k \oplus p, x) = \chi_{\mathrm{band}}(k, x) \chi_{\mathrm{band}}(p, x).
    \end{equation}
    \item \textbf{Complex Conjugation:} If $K = \mathbb{C}$ and $|\zeta| = 1$, then $\overline{\chi(k, x)} = \chi(-k, x)$.
\end{enumerate}
\end{theorem}

\begin{proof}
Items 1--3 follow directly from Lemma~\ref{lem:a01_rep_independence} and the index laws $\zeta^{k(x+y)} = \zeta^{kx}\zeta^{ky}$.
For item 4: by Theorem~\ref{thm:wrap_lemma}, $k \oplus p = k + p - w L$ where $w \in \{-1, 0, 1\}$.
By Lemma~\ref{lem:a01_rep_independence}:
\[
\chi_{\mathrm{band}}(k \oplus p, x) = \zeta^{(k + p - w L) x_{\mathrm{val}}} = \zeta^{(k+p) x_{\mathrm{val}}} = \zeta^{k x_{\mathrm{val}}} \zeta^{p x_{\mathrm{val}}} = \chi_{\mathrm{band}}(k, x) \chi_{\mathrm{band}}(p, x).
\]
For item 5: if $|\zeta| = 1$, then $\overline{\zeta} = \zeta^{-1}$, so $\overline{\zeta^{kx}} = (\zeta^{-1})^{kx} = \zeta^{-kx} = \chi(-k, x)$.
\end{proof}

\section{Geometric Sum Cancellation and Dual Orthogonality}

\begin{lemma}[Non-Trivial Character Cancellation]
\label{lem:nontrivial_character_sum}
Let $K$ be a field, and let $\zeta \in K$ be a primitive $L$-th root of unity. For any residue $k \in \mathbb{Z}/L\mathbb{Z}$:
\begin{equation}
\sum_{x \in \mathbb{Z}/L\mathbb{Z}} \chi_{\mathrm{res}}(k, x) = \begin{cases} L \cdot 1_K, & \text{if } k = 0, \\ 0, & \text{if } k \neq 0. \end{cases}
\end{equation}
\end{lemma}

\begin{proof}
If $k = 0$, then $\chi_{\mathrm{res}}(0, x) = \zeta^0 = 1$ for all $x$. The sum of $L$ copies of $1_K$ is $L \cdot 1_K$.

Now assume $k \neq 0$. Since $\zeta$ is a primitive $L$-th root of unity and $k \not\equiv 0 \pmod L$, $\zeta^{k_{\mathrm{val}}} \neq 1$.
Consider the finite sum $S = \sum_{x=0}^{L-1} (\zeta^{k_{\mathrm{val}}})^x$. Multiplying by $(\zeta^{k_{\mathrm{val}}} - 1)$:
\[
(\zeta^{k_{\mathrm{val}}} - 1) S = \sum_{x=0}^{L-1} (\zeta^{k_{\mathrm{val}}})^{x+1} - \sum_{x=0}^{L-1} (\zeta^{k_{\mathrm{val}}})^x = (\zeta^{k_{\mathrm{val}}})^L - 1 = (\zeta^L)^{k_{\mathrm{val}}} - 1 = 1^{k_{\mathrm{val}}} - 1 = 0.
\]
Because $K$ is a field and $\zeta^{k_{\mathrm{val}}} - 1 \neq 0$, we can divide by $(\zeta^{k_{\mathrm{val}}} - 1)$ to conclude $S = 0$.
\end{proof}

\begin{theorem}[Spatial and Momentum Orthogonality]
\label{thm:app_orthogonality}
Let $K$ be a field, $\zeta \in K$ a primitive $L$-th root of unity, and $h = L/2 \ge 1$.
\begin{enumerate}
    \item \textbf{Spatial Orthogonality on $\Lambda^*$:} For any two momentum modes $k, p \in \Lambda^*$:
    \begin{equation}
    \sum_{x \in \Lambda} \chi_{\mathbb{Z}}(\zeta, k - p, x_{\mathrm{val}}) = L \delta_{k, p} \cdot 1_K.
    \end{equation}
    \item \textbf{Momentum Orthogonality on $\Lambda$:} For any two lattice sites $x, y \in \Lambda$:
    \begin{equation}
    \sum_{k \in \Lambda^*} \chi_{\mathbb{Z}}(\zeta, k, (x - y)_{\mathrm{val}}) = L \delta_{x, y} \cdot 1_K.
    \end{equation}
\end{enumerate}
\end{theorem}

\begin{proof}
\begin{enumerate}
    \item By Theorem~\ref{thm:character_homomorphisms}, $\chi_{\mathbb{Z}}(\zeta, k - p, x_{\mathrm{val}}) = \chi_{\mathrm{res}}(\pi(k) - \pi(p), x)$.
    By Lemma~\ref{lem:nontrivial_character_sum}, the sum over $x$ equals $L \cdot 1_K$ if $\pi(k) = \pi(p)$, and $0$ otherwise.
    Since the projection $\pi: \Lambda^* \to \mathbb{Z}/L\mathbb{Z}$ is a bijection with inverse $r$ (Lemma~\ref{lem:section_identities}), $\pi(k) = \pi(p)$ if and only if $k = p$.
    \item Using the bijection $\pi: \Lambda^* \to \mathbb{Z}/L\mathbb{Z}$, summing $k$ over $\Lambda^*$ is identical to summing the residue $\kappa$ over $\mathbb{Z}/L\mathbb{Z}$.
    Applying Lemma~\ref{lem:nontrivial_character_sum} with role of $k$ and $x$ exchanged yields $L \delta_{x, y} \cdot 1_K$.
\end{enumerate}
\end{proof}

\section{Unscaled Transforms and Two-Sided Inversion}

\begin{definition}[Unscaled Forward and Backward Transforms]
Let $K^\Lambda$ and $K^{\Lambda^*}$ denote the $K$-vector spaces of functions on the lattice and band, respectively.
\begin{enumerate}
    \item The \emph{unscaled analysis map} $S: K^\Lambda \to K^{\Lambda^*}$ is defined by:
    \begin{equation}
    (S f)(k) \coloneqq \sum_{x \in \Lambda} f(x) \chi_{\mathbb{Z}}(\zeta, -k, x_{\mathrm{val}}).
    \end{equation}
    \item The \emph{unscaled synthesis map} $T: K^{\Lambda^*} \to K^\Lambda$ is defined by:
    \begin{equation}
    (T g)(x) \coloneqq \sum_{k \in \Lambda^*} g(k) \chi_{\mathbb{Z}}(\zeta, k, x_{\mathrm{val}}).
    \end{equation}
\end{enumerate}
\end{definition}

\begin{theorem}[Two-Sided Unscaled Inversion]
\label{thm:a01_unscaled_inversion}
The linear maps $S$ and $T$ satisfy:
\begin{equation}
T \circ S = L \cdot I_{K^\Lambda}, \qquad S \circ T = L \cdot I_{K^{\Lambda^*}}.
\end{equation}
Whenever $L \cdot 1_K \neq 0$ in $K$, $S$ is a linear isomorphism with inverse $S^{-1} = L^{-1} T$.
\end{theorem}

\begin{proof}
Let $f \in K^\Lambda$ and evaluate $(T(Sf))(y)$ for $y \in \Lambda$:
\begin{align*}
(T(Sf))(y) &= \sum_{k \in \Lambda^*} (Sf)(k) \chi(k, y) = \sum_{k \in \Lambda^*} \left( \sum_{x \in \Lambda} f(x) \chi(-k, x) \right) \chi(k, y) \\
&= \sum_{x \in \Lambda} f(x) \left( \sum_{k \in \Lambda^*} \chi(k, y - x) \right).
\end{align*}
By Theorem~\ref{thm:app_orthogonality}(2), the inner sum is $L \delta_{x, y} \cdot 1_K$.
Therefore:
\[
(T(Sf))(y) = \sum_{x \in \Lambda} f(x) (L \delta_{x, y} \cdot 1_K) = L f(y).
\]
Hence $T \circ S = L I$. A completely symmetric calculation exchanging the order of summation and using Theorem~\ref{thm:app_orthogonality}(1) proves $S \circ T = L I$.
\end{proof}

\section{Exact Dispersion Algebra in Arbitrary Fields}

In physical lattices, the nearest-neighbor discrete Laplacian is given by $(\Delta f)(x) = f(x+1) + f(x-1) - 2 f(x)$.
Its eigenvalue on plane-wave modes is determined purely algebraically without trigonometry.

\begin{lemma}[Exact Second-Difference Eigenvalue]
\label{lem:laplacian_algebraic}
Let $K$ be a field and $\zeta \in K^\times$. For any integer $k \in \mathbb{Z}$:
\begin{equation}
(\zeta^k - 1)(1 - \zeta^{-k}) = \zeta^k + \zeta^{-k} - 2.
\end{equation}
\end{lemma}

\begin{proof}
Expanding the left-hand side:
\[
(\zeta^k - 1)(1 - \zeta^{-k}) = \zeta^k \cdot 1 - \zeta^k \zeta^{-k} - 1 \cdot 1 + 1 \cdot \zeta^{-k} = \zeta^k - 1 - 1 + \zeta^{-k} = \zeta^k + \zeta^{-k} - 2.
\]
\end{proof}

\begin{proposition}[Connection to Standard Trigonometric Dispersion]
Specializing to $K = \mathbb{C}$ with the canonical root $\zeta = e^{2\pi i / L}$:
\begin{equation}
\zeta^k + \zeta^{-k} - 2 = - 4 \sin^2\left(\frac{\pi k}{L}\right).
\end{equation}
\end{proposition}

\begin{proof}
Using Euler's formula:
\[
\zeta^k + \zeta^{-k} - 2 = 2 \cos\left(\frac{2\pi k}{L}\right) - 2 = -2 \left(1 - \cos\left(\frac{2\pi k}{L}\right)\right).
\]
Using the standard double-angle identity $1 - \cos(2\theta) = 2 \sin^2(\theta)$ with $\theta = \frac{\pi k}{L}$ yields:
\[
-2 \left(2 \sin^2\left(\frac{\pi k}{L}\right)\right) = - 4 \sin^2\left(\frac{\pi k}{L}\right).
\]
\end{proof}

\begin{remark}[Formalization Note]
In the verified Lean~4 core (Appendix~A01), all character homomorphisms, the representative-independence theorem, the geometric sum cancellation, and the unscaled operator identities are proved without noncomputable choice over generic fields. The physical normalization $a = 1/\sqrt{L}$ is introduced only in the complex Hilbert tier to construct the unitary isometry $U: \ell^2(\Lambda) \to \ell^2(\Lambda^*)$.
\end{remark}
